%%%%%%%%% ABSTRACT
\begin{abstract}
When an inter-symbol interference (ISI) channel changes with time, the optimum equalizer becomes a moving target. This project asks how fast that target may move before a fixed-point least mean squares (LMS) equalizer, intended for an FPGA, can no longer follow it. A baseband BPSK link with a deliberately time-varying channel feeds three receivers on one impaired stream: a direct detector, an LMS equalizer trained once and frozen, and a continuously adapting LMS equalizer. Bit-error rate, mean-square error and reacquisition time are measured against the variation rate, the step size and the coefficient word length. This mid-term report presents the further literature review, the published results replicated, and the software results of phases P1--P6.
\nb{All results are software: floating-point simulation (\lvl{L1}), bit-true fixed-point simulation (\lvl{L2}) and, as early progress, RTL simulation (\lvl{L3}); synthesis (\lvl{L4}) and physical FPGA tests (\lvl{L5}) have not been done. At \lvl{L1} (100 seeds) a tracking boundary is located for every step size: for $\mu\le2^{-7.5}$, on both channels at 14 and 16\,dB, it follows $S^\star_{\rm MSE}\approx0.11/\mu$, so halving the step size requires a channel that changes about half as often. A first-order theory model with no fitted constant reproduces this boundary, at the resolution of the tested rate grid, in all 40 such cases, and explains the measured optimum-step exponent: about 0.52 on the primary channel, consistent with 1/2 within its systematic uncertainty, and 0.58--0.59 on the control channel, between 1/2 and 2/3. The textbook misadjustment formula is found to be 1.47 times too small on an ISI channel, and a correlated-error correction is derived. At \lvl{L2} (16\,dB, 20 seeds, round-to-nearest updates, 12-bit multiplier), a 16-bit coefficient accumulator preserves the MSE tracking boundary of the 26-bit baseline at all seven power-of-two step sizes, whereas the multiplier width affects only static accuracy; this is a software-supported design recommendation, not a measured hardware result. Of seven published results rebuilt, six are reproduced (two with stated exceptions) and one documented mismatch is explained; MATLAB confirms the Python LMS model independently.} The recorded-audio and image payloads and decision-directed operation are also reported.
\end{abstract}

%%%%%%%%% BODY
\section{Introduction}
\label{sec:intro}

A channel with memory spreads each transmitted symbol over several symbol intervals. This inter-symbol interference does not fall as transmit power rises, so a fixed-threshold detector keeps an error floor that adaptive equalization removes by learning an approximate inverse of the channel \cite{lucky1965,qureshi1985}. For a static channel the optimum equalizer $\mathbf w^\star$ is a fixed point and LMS \cite{widrowhoff1960} converges to it. When the channel moves --- through motion, changing reflections or propagation conditions --- the optimum becomes $\mathbf w^\star[n]$, and the question changes from \emph{does the filter converge} to \emph{can the filter adapt faster than the channel changes} \cite{widrow1976,macchi1995}. The variation rate used here is a laboratory analogue of the normalised Doppler $f_DT_s$ that governs equalizer tracking in mobile, underwater-acoustic and HF links \cite{jakes1974}.

The preliminary report (PPR) posed the research question, the channel model, the three-receiver experiment, the fixed-point and FPGA architecture for the Real Digital Boolean Board (Spartan-7 XC7S50) \cite{realdigitalboard}, and a development plan in which phases P1--P6 (software) form the mid-term target and P7--P10 (RTL, synthesis, board) the final target. This report does not repeat that material. It gives the methodology needed to read the results (\cref{sec:method}) and then reports what has been examined, replicated, measured and explained since the PPR.

\noindent\textbf{Progress since the PPR.} All software phases P1--P6 are complete and P7--P8 have started:
\begin{itemize}\itemsep0pt
\item a further, phase-by-phase literature review and an audit of the PPR references (\cref{sec:literature});
\item seven published results rebuilt and compared with their sources (\cref{sec:replication});
\item P1--P4: validation, abrupt transitions, the tracking boundary against step size and SNR, the optimum step size, coefficient trajectories and decision-directed operation (\cref{sec:results14});
\item P5--P6: the recorded-voice and image payloads and the bit-true fixed-point word-length study, plus a first RTL chain (\cref{sec:results56});
\item \nb{a theory model that predicts the measured boundary and explains the optimum-step exponent (\cref{sec:theory}), and a correction of the misadjustment formula (\cref{sec:misadj})}.
\end{itemize}
Conclusions, limitations and the remaining hardware work close the report (\cref{sec:conclusion}); every change to a PPR statement is listed in \cref{app:register}.

\noindent\textbf{Reading this report.} \emph{Black} text is material stated, planned or cited in the PPR, including PPR plans that have now been carried out. \nb{\emph{Blue} text is new: new ideas and methods, newly reviewed literature, and the mathematics, calculations and numerical results worked out since the PPR.} Every number carries its validation level: \lvl{L1} floating-point simulation, \lvl{L2} bit-true fixed-point simulation, \lvl{L3} RTL simulation; \lvl{L4} (synthesis) and \lvl{L5} (board) results do not exist yet. Unless stated otherwise, adaptation is \emph{supervised}: R3 uses the known transmitted symbol as its reference (\cref{sec:receivers}).
